\title{Group Sequence Policy Optimization}

\begin{abstract}
We present Group Sequence Policy Optimization (GSPO), a robust and effective reinforcement learning method for training large language models. Distinct from prior approaches that use token-wise importance ratios, GSPO utilizes sequence-level likelihood to define importance ratios, and applies clipping, rewarding, and optimization at the sequence level. Our results show that GSPO provides notable improvements in training speed and performance over the GRPO algorithm, delivers significant stabilization in Mixture-of-Experts (MoE) RL training, and offers opportunities to streamline reinforcement learning infrastructure design. These advantages have played a key role in the significant advancements observed in the latest Qwen3 models.
\end{abstract}

\section{Introduction}
\label{sec:intro}

Reinforcement learning (RL) has become a central framework for advancing the scalability of language models \citep{o1,dpsk-r1,qwq32b,qwen3}. Leveraging large-scale RL allows language models to acquire proficiency in complex tasks—including competitive mathematics and programming—by enabling extended and more intricate chains of reasoning.

To harness the benefits of increased computational resources for RL, it is critical to ensure that training remains stable and reliable. Nevertheless, the latest RL approaches, such as GRPO \citep{grpo}, face pronounced stability challenges when applied to extremely large language models. These difficulties frequently manifest as catastrophic and unrecoverable collapse of the model \citep{qwen3, minimax-m1}, posing a significant obstacle to advancing the performance of language models through further RL-based optimization.

In this work, we demonstrate that the root cause of instability in GRPO originates from a fundamental misuse and consequent invalidation of importance sampling weights within its algorithmic structure. This flaw leads to substantial training variance, which intensifies as response sequences grow longer, and the problem is exacerbated by the use of a clipping mechanism—eventually resulting in model collapse.

To overcome these fundamental issues, we introduce \textbf{Group Sequence Policy Optimization (GSPO)}, a novel RL algorithm designed for training large language models. The principal advancement of GSPO is its rigorously defined importance ratio, which is based on sequence likelihood \citep{click}, thereby adhering to the foundational concepts of importance sampling. Furthermore, GSPO calculates normalized rewards as the advantages derived from multiple responses to a single query, facilitating coherent alignment between sequence-level rewards and the optimization process.

Our experimental results highlight that GSPO substantially outperforms GRPO with respect to training robustness, efficiency, and overall effectiveness. Notably, GSPO intrinsically addresses the instability issues commonly encountered when applying RL methods to large Mixture-of-Experts (MoE) models, thereby removing the reliance on intricate stabilization mechanisms. This capability also suggests a pathway to streamline RL systems. As a result, GSPO has been instrumental in achieving considerable performance gains in the most recent iterations of the Qwen3 models. We anticipate that GSPO will serve as a solid and scalable basis for future progress in large-scale RL training for language models.

\section{Preliminaries}

\paragraph{Notation}
Throughout this work, we represent an autoregressive language model with parameters $\theta$ as a policy $\pi_\theta$.
The symbol $x$ signifies a query, and $\mathcal{D}$ refers to the collection of queries. For a particular response $y$ to query $x$, the probability assigned by the policy $\pi_\theta$ is expressed as $\pi_\theta (y | x)=\prod_{t=1}^{|y|} \pi_\theta (y_t | x, y_{<t} )$, where $|y|$ indicates the total number of tokens in the response $y$. Any pair $(x, y)$ consisting of a query and its response can be evaluated by a verifier $r$, which assigns a reward $r(x, y)$ ranging within $[ 0, 1 ]$.

\paragraph{Proximal Policy Optimization (PPO)} Using samples generated from the old policy $\pi_{\theta_\text{old}}$, PPO \citep{ppo} constrains the policy update within a proximal region of the old policy through the clipping mechanism. Specifically, PPO employs the following objective for policy optimization (we omit the KL regularization term hereinafter for brevity, as it is not the focus of this paper):

\begin{align}
\mathcal{J}_\text{PPO}(\theta) = \mathbb{E}_{ x \sim \mathcal{D},\, y \sim \pi_{\theta_\text{old}}( \cdot | x) }
\left[ \frac{1}{|y|} \sum_{t=1}^{|y|} 
\min \left( w_{t}(\theta) \widehat{A}_{t},  \, \mathrm{clip} \left( w_{t}(\theta), 1 - {\varepsilon}, 1 + {\varepsilon}\right) \widehat{A}_{t} \right)
\right],
\end{align}

The importance ratio for token $y_t$ is expressed as $
w_{t}(\theta) = \frac{ \pi_{\theta} (y_{t} | x, y_{<t}) }{ \pi_{\theta_\text{old}} (y_{t} | x,y_{<t})} $, with the advantage $\widehat{A}_{t}$ of $y_t$ obtained through a separate value model, and $\varepsilon$ representing the range used to clip importance ratios.

In practical applications, PPO faces a fundamental difficulty owing to its substantial dependence on the value model. Typically, the value model is comparable in scale to the policy model, thereby increasing both the computational and memory overhead. Moreover, the efficacy of PPO is tightly linked to the precision of its value estimates. Achieving a dependable value model is nontrivial in itself, and scaling such reliability to scenarios requiring longer outputs and tackling more sophisticated tasks is even more problematic.

\paragraph{Group Relative Policy Optimization (GRPO)} GRPO \citep{grpo} bypasses the need for the value model by computing the relative advantage of each response within a group of responses to the same query. Specifically, GRPO optimizes the following objective:

\begin{align}
\label{equ:grpo}
\mathcal{J}_\text{GRPO}(\theta) = \mathbb{E}_{ x \sim \mathcal{D},\, \{y_i\}_{i=1}^G \sim \pi_{\theta_\text{old}}( \cdot | x) }
\left[ \frac{1}{G} \sum_{i=1}^{G} \frac{1}{|y_i|} \sum_{t=1}^{|y_i|} 
\min \left( w_{i,t}(\theta) \widehat{A}_{i,t},  \, \mathrm{clip} \left( w_{i,t}(\theta), 1 - {\varepsilon}, 1 + {\varepsilon}\right) \widehat{A}_{i,t} \right)
\right],
\end{align}

where $G$ is the number of generated responses to each query $x$ (i.e., the group size), and the importance ratio $w_{i,t}(\theta)$ and advantage $\widehat{A}_{i,t}$ of token $y_{i,t}$ are:

\begin{align}
    w_{i,t}(\theta)=\frac{ \pi_{\theta} (y_{i,t} | x, y_{i,<t}) }{ \pi_{\theta_\text{old}} (y_{i,t} | x,y_{i,<t})},\quad\ 
    \widehat{A}_{i,t} = \widehat{A}_{i} = \frac{r(x, y_i) - \mathrm{mean} \left( \{ r(x, y_i) \}_{i=1}^G \right) }{ \mathrm{std} \left( \{ r(x, y_i) \}_{i=1}^G \right) },
\end{align}

respectively, where all the tokens in $y_i$ share the same advantage as $\widehat{A}_{i}$.

\section{Motivation}
\label{sec:motivation}

As model architectures expand in terms of parameter count, incorporate sparsity (such as in Mixture-of-Experts frameworks), and produce longer outputs, efficient hardware usage during reinforcement learning requires scaling up the rollout batch size. To enhance sample efficiency, it is customary to subdivide these sizable rollout batches into several smaller mini-batches for use in gradient computation. This approach, however, leads to an off-policy training regime: the outputs $y$ are generated by a previous policy $\pi_{\theta_\text{old}}$ and not by the currently updated policy $\pi_\theta$. Such off-policy scenarios highlight why clipping methods, as implemented in PPO and GRPO, are crucial—they serve to restrict the impact of samples that are significantly misaligned with the ongoing policy, thereby maintaining reliable gradient estimates.

While mechanisms like clipping aim to manage this off-policy discrepancy, we identify a more fundamental issue in GRPO: \textit{its objective is ill-posed}. This problem becomes particularly acute when training large models on long-response tasks, leading to catastrophic model collapse. The ill-posed nature of the GRPO objective stems from a misapplication of importance sampling weights. The principle of importance sampling is to estimate the expectation of a function $f$ under a target distribution $\pi_\text{tar}$ by re-weighting samples drawn from a behavior distribution $\pi_\text{beh}$:

\begin{align}
\mathbb{E}_{z \sim \pi_\text{tar}} \left[ f(z) \right] = \mathbb{E}_{z \sim \pi_\text{beh}} \left[ \frac{ \pi_\text{tar} (z) }{ \pi_\text{beh} (z)} \, f(z) \right].
\end{align}

The effectiveness of this approach fundamentally depends on aggregating results over many samples ($N \gg 1$) drawn from the behavior distribution $\pi_\text{beh}$, allowing the importance weight $\frac{ \pi_\text{tar} (z) }{ \pi_\text{beh} (z)}$ to properly address disparities between the behavior and target distributions.

By contrast, GRPO employs the importance weight $\frac{ \pi_{\theta} (y_{i,t} | x, y_{i,<t}) }{ \pi_{\theta_\text{old}} (y_{i,t} | x,y_{i,<t})}$ at each token timestep $t$. This weight is computed based solely on a single realization $y_{i,t}$ sampled from the conditional next-token distribution $\pi_{\theta_\text{old}}( \cdot | x, y_{i, <t})$, preventing it from functioning as an appropriate distributional correction. Rather, it injects considerable variance into the gradient estimates during training—a phenomenon that intensifies as sequence lengths increase and is further aggravated by the presence of the clipping procedure. Empirically, we find that this high-variance behavior often triggers model collapse, after which recovery is extremely difficult. Attempts to continue training from previous checkpoints, including careful adjustment of hyperparameters such as clipping thresholds, extending generated sequence length, or altering RL sampling strategies, do not reverse the collapse once it occurs.

These findings highlight a fundamental flaw in GRPO's construction: the mismatch between the granularity of the importance weighting and the reward. In particular, this suggests the essential principle that \textbf{the scale of the optimization target must correspond with the scale at which rewards are issued}. Since the reward signals are attributed to complete sequences, applying off-policy correction on a per-token basis leads to significant issues. Consequently, our analysis directs attention away from token-level objectives and advocates for adopting importance weighting and performing optimization at the \textit{sequence level}.

\section{Algorithm}

\subsection{GSPO: Group Sequence Policy Optimization}

Although the use of token-level importance weights $\frac{ \pi_{\theta} (y_{i,t} | x, y_{i,<t}) }{ \pi_{\theta_\text{old}} (y_{i,t} | x, y_{i,<t})}$ presents challenges within GRPO, we note that, for language generation tasks, the \textit{sequence-level} importance weight $\frac{ \pi_{\theta} (y | x) }{ \pi_{\theta_\text{old}} (y | x)}$ possesses a well-defined theoretical interpretation: it quantifies the extent to which a response $y$ drawn from $\pi_{\theta_\text{old}} (\cdot | x)$ diverges from $\pi_{\theta} (\cdot | x)$. This perspective resonates with sequence-level rewards and also provides an interpretable basis for understanding the effects of the clipping procedure.

Based on this straightforward observation, we propose the \textbf{Group Sequence Policy Optimization (GSPO)} algorithm. GSPO employs the following sequence-level optimization objective:

\begin{align}
\mathcal{J}_\text{GSPO} (\theta) =
\mathbb{E}_{ x \sim \mathcal{D},\, \{y_i\}_{i=1}^G \sim \pi_{\theta_\text{old}}( \cdot | x) }
\left[ 
\frac{1}{G} \sum_{i=1}^{G}
\min \left( s_{i}(\theta)  \widehat{A}_{i},  \, \mathrm{clip} \left( s_{i}(\theta), 1 - {\varepsilon}, 1 + {\varepsilon} \right) \widehat{A}_{i} \right) 
\right],
\label{equ:gspo}
\end{align}

where we adopt the group-based advantage estimation:

\begin{align}
\widehat{A}_{i} = \frac{r(x, y_i) - \mathrm{mean} \left( \{ r(x, y_i) \}_{i=1}^G \right) }{ \mathrm{std} \left( \{ r(x, y_i) \}_{i=1}^G \right) },
\end{align}

and define the importance ratio $s_{i}(\theta)$ based on sequence likelihood \citep{click}:

\begin{align}
s_{i}(\theta) = \left( \frac{ \pi_{\theta} (y_i | x) }{ \pi_{\theta_\text{old}} (y_i | x)} \right)^{\frac{1}{|y_i|}}
=
\exp \left( \frac{1}{|y_i|} \sum_{t=1}^{|y_i|} \log \frac{ \pi_{\theta} (y_{i,t} | x, y_{i,<t}) }{ \pi_{\theta_\text{old}} (y_{i,t} | x,y_{i,<t})} \right).
\end{align}

Consequently, GSPO utilizes response-level clipping, rather than token-level clipping, to filter out highly ``off-policy'' samples from gradient calculation, thereby aligning with both sequence-level reward allocation and optimization objectives. Importantly, we implement length normalization on $s_{i}(\theta)$ to mitigate variance and to keep $s_{i}(\theta)$ confined to a consistent numerical interval. In the absence of length normalization, variations in the likelihood of certain tokens could cause significant volatility in the sequence-level importance ratio, necessitating different clipping thresholds for responses of varying lengths. Additionally, it is worth noting that the clipping thresholds in GSPO and prior methods such as GRPO generally differ significantly in magnitude, reflecting the disparate formalizations of importance ratios.

\subsection{Gradient Analysis}
\label{subsec:gradient}

We can derive the gradient of the GSPO objective as follows (clipping is omitted for brevity):

\begin{align}
\nabla_{\theta} \mathcal{J}_\text{GSPO} (\theta)
=&\ 
\nabla_{\theta} \mathbb{E}_{ x \sim \mathcal{D},\, \{y_i\}_{i=1}^G \sim \pi_{\theta_\text{old}}( \cdot | x) }
\left[ 
\frac{1}{G} \sum_{i=1}^{G} s_{i}(\theta) \widehat{A}_{i}
\right] \\
= &\ 
\mathbb{E}_{ x \sim \mathcal{D},\, \{y_i\}_{i=1}^G \sim \pi_{\theta_\text{old}}( \cdot | x) }
\left[ 
\frac{1}{G} \sum_{i=1}^{G}
s_{i}(\theta) \widehat{A}_{i}
\cdot \nabla_{\theta} \log s_{i}(\theta)
\right] \\
=&\ 
\mathbb{E}_{ x \sim \mathcal{D},\, \{y_i\}_{i=1}^G \sim \pi_{\theta_\text{old}}( \cdot | x) }
\left[ 
\frac{1}{G} \sum_{i=1}^{G} 
\left( \frac{ \pi_{\theta} (y_i | x) }{ \pi_{\theta_\text{old}} (y_i | x)} \right)^{\frac{1}{|y_i|}} \widehat{A}_{i} 
\cdot \frac{1}{|y_i|} \sum_{t=1}^{|y_i|} \nabla_{\theta} \log \pi_{\theta} (y_{i,t} | x, y_{i,<t}) 
\right].
\label{equ:gspo_grad}
\end{align}

For comparison, the gradient of the GRPO objective is as follows (note that $\widehat{A}_{i,t} = \widehat{A}_{i}$):

\begin{align}
\nabla_{\theta} \mathcal{J}_\text{GRPO} (\theta) 
=&\ 
\nabla_{\theta} \mathbb{E}_{ x \sim \mathcal{D},\, \{y_i\}_{i=1}^G \sim \pi_{\theta_\text{old}}( \cdot | x) }
\left[ \frac{1}{G} \sum_{i=1}^{G} \frac{1}{|y_i|} \sum_{t=1}^{|y_i|} 
w_{i,t}(\theta) \widehat{A}_{i,t}
\right] \\
=&\
\mathbb{E}_{ x \sim \mathcal{D},\, \{y_i\}_{i=1}^G \sim \pi_{\theta_\text{old}}( \cdot | x) }
\left[ \frac{1}{G} \sum_{i=1}^{G} \widehat{A}_{i} 
\cdot \frac{1}{|y_i|} \sum_{t=1}^{|y_i|} 
\frac{ \pi_{\theta} (y_{i,t} | x, y_{i,<t}) }{ \pi_{\theta_\text{old}} (y_{i,t} | x,y_{i,<t})} 
\nabla_{\theta} \log \pi_{\theta} (y_{i,t} | x, y_{i,<t})  
\right].
\end{align}

The core difference between GSPO and GRPO is found in \textit{the method by which each algorithm assigns weights to the gradients of token log likelihoods}. GRPO utilizes an ``importance weight'' for each token, specifically $\frac{ \pi_{\theta} (y_{i,t} | x, y_{i,<t}) }{ \pi_{\theta_\text{old}} (y_{i,t} | x,y_{i,<t})}$, which varies from token to token. These non-uniform weights, which can span the intervals $(0, 1+\varepsilon]$ (when $\widehat{A}_{i} >0$) or $[1-\varepsilon, +\infty)$ (when $\widehat{A}_{i} < 0$), can significantly influence the learning process, with cumulative effects potentially resulting in erratic behavior during training. By contrast, GSPO applies a uniform weighting scheme across all tokens within a response, thereby removing the instability associated with GRPO’s variable weighting.

\subsection{GSPO-token: A Token-level Objective Variant}

In scenarios like multi-turn RL, we may desire a finer-grained advantage adjustment than the sequence level. To this end, we introduce a token-level objective variant of GSPO, namely \textbf{GSPO-token}, to allow token-wise advantage customization:

\begin{align}
\mathcal{J}_\text{GSPO-token}(\theta)
=
\mathbb{E}_{ x \sim \mathcal{D},\, \{y_i\}_{i=1}^G \sim \pi_{\theta_\text{old}}( \cdot | x) }
\left[ 
\frac{1}{G} \sum_{i=1}^{G} \frac{1}{|y_i|} \sum_{t=1}^{|y_i|} 
\min \left( s_{i,t}(\theta) \widehat{A}_{i,t},  \, \mathrm{clip} \left( s_{i,t}(\theta), 1 - {\varepsilon}, 1 + {\varepsilon} \right) \widehat{A}_{i,t} \right) 
\right],
\label{equ:gspo-token}
\end{align}

where

\begin{align}
s_{i,t}(\theta) 
= \mathrm{sg} \left[ s_{i}(\theta) \right]  \cdot \frac{ \pi_{\theta} (y_{i,t} | x, y_{i,<t}) }{ \mathrm{sg} \left[ \pi_{\theta} (y_{i,t} | x, y_{i,<t}) \right] },
\end{align}

and $\mathrm{sg}[\cdot]$ denotes only taking the numerical value but stopping the gradient, corresponding to the \texttt{detach} operation in PyTorch. The gradient of GSPO-token can be derived as:

\begin{align}
\nabla_{\theta} \mathcal{J}_\text{GSPO-token}(\theta)
=&\ 
\nabla_{\theta} \mathbb{E}_{ x \sim \mathcal{D},\, \{y_i\}_{i=1}^G \sim \pi_{\theta_\text{old}}( \cdot | x) }
\left[ 
\frac{1}{G} \sum_{i=1}^{G}
\frac{1}{|y_i|} \sum_{t=1}^{|y_i|} 
s_{i,t}(\theta) \widehat{A}_{i,t}
\right] \\
=&\ 
\mathbb{E}_{ x \sim \mathcal{D},\, \{y_i\}_{i=1}^G \sim \pi_{\theta_\text{old}}( \cdot | x) }
\left[ 
\frac{1}{G} \sum_{i=1}^{G} s_i (\theta) \cdot \frac{1}{|y_i|} \sum_{t=1}^{|y_i|} 
\widehat{A}_{i,t} \frac{ \nabla_{\theta} \pi_{\theta} (y_{i,t} | x, y_{i,<t}) }{ \pi_{\theta} (y_{i,t} | x, y_{i,<t}) }
\right] \\
=&\ 
\mathbb{E}_{ x \sim \mathcal{D},\, \{y_i\}_{i=1}^G \sim \pi_{\theta_\text{old}}( \cdot | x) }
\left[ 
\frac{1}{G} \sum_{i=1}^{G}  \left( \frac{ \pi_{\theta} (y_i | x) }{ \pi_{\theta_\text{old}} (y_i | x)} \right)^{\frac{1}{|y_i|}}
\cdot \frac{1}{|y_i|} \sum_{t=1}^{|y_i|}
\widehat{A}_{i,t} \nabla_{\theta} \log \pi_{\theta} (y_{i,t} | x, y_{i,<t})  
\right].
\label{equ:gspo-token_grad}
\end{align}

It is important to observe that the expression $\frac{ \pi_{\theta} (y_{i,t} | x, y_{i,<t}) }{ \mathrm{sg} \left[ \pi_{\theta} (y_{i,t} | x, y_{i,<t}) \right] }$ evaluates numerically to 1, which implies that $s_{i,t}(\theta)$ is numerically indistinguishable from $s_{i}(\theta)$. If we analyze Equation~\eqref{equ:gspo} alongside~\eqref{equ:gspo-token}, as well as Equation~\eqref{equ:gspo_grad} with~\eqref{equ:gspo-token_grad}, we find that GSPO-token and GSPO yield identical values for the objective function, clipping mechanism, and theoretical gradient, provided that the advantage for each token in the generated output $y_i$ is held constant across the sequence (i.e., $\widehat{A}_{i,t} = \widehat{A}_{i}$). In contrast, GSPO-token enables finer granularity by permitting independent adjustment of the advantage at the token level.

\section{Experiments and Discussion}

\subsection{Empirical Results}

We conduct experiments using a cold-start model initialized from Qwen3-30B-A3B-Base and provide both the training reward trajectories and evaluation metrics on three benchmarks: AIME'24 (mean Pass@1 across 32 samples), LiveCodeBench (202410-202502, mean Pass@1 across 8 samples), and CodeForces (Elo Rating). Throughout the reinforcement learning phase, each set of rollout data is divided into four mini-batches to facilitate gradient computations. For GSPO, we specify the clipping intervals in Equation~\eqref{equ:gspo} as 3e-4 (left) and 4e-4 (right). Our baseline is GRPO, for which the clipping bounds in Equation~\eqref{equ:grpo} are set to 0.2 and 0.27, respectively, after thorough hyperparameter tuning to maintain comparability. It is important to highlight that GRPO requires implementation of the Routing Replay training method to enable reliable convergence for MoE RL, a topic elaborated upon in \S~\ref{subsec:routing_replay}. In contrast, \textit{GSPO eliminates the necessity for this approach}.

As illustrated in Figure~\ref{fig:results}, GSPO enables consistently stable progression during training. Our findings indicate that \textbf{GSPO supports ongoing enhancement in performance by scaling up training compute, periodically refreshing the query set, and lengthening generation output}. In addition, GSPO exhibits greater training efficiency compared to GRPO, attaining higher accuracy and better results on benchmarks given equal training compute and query usage. We also report the successful integration of GSPO into RL-based training for the latest Qwen3 models, offering strong validation of GSPO’s effectiveness in leveraging RL scaling for large language models.

\subsection{Curious Observation on Clipping Fractions}

An important feature that sets GSPO apart from GRPO is its method of clipping responses at the sequence level, rather than at the token level. As illustrated in Figure~\ref{fig:clipping}, the proportion of clipped tokens in GSPO is smaller by approximately two orders of magnitude compared to GRPO, even when the clipping ranges are modified; this does not diminish the observed difference in magnitude. Intriguingly, despite GSPO discarding substantially more tokens and thus utilizing a smaller portion of data for optimization or gradient calculation, it demonstrates greater training efficiency over GRPO. This surprising result—that training with a markedly higher fraction of clipped tokens results in improved efficiency—suggests that the token-wise gradient estimates provided by GRPO suffer from greater noise and are less optimal for exploiting samples. GSPO’s strategy, which relies on sequence-level gradients, delivers a stronger and more consistent learning signal.

\subsection{Benefit of GSPO for MoE Training}
\label{subsec:routing_replay}

\paragraph{Background}
MoE models, with their characteristic sparse activation patterns, present distinct stability issues in RL training not observed in dense models. Specifically, we observed that when utilizing the GRPO algorithm, MoE models are susceptible to \textit{expert-activation volatility}, which often obstructs the proper convergence of RL. In detail, even after only a few rounds of gradient updates, there can be substantial shifts in the set of experts triggered for the same response. Taking the 48-layer Qwen3-30B-A3B-Base model as an example, post RL gradient update, approximately 10\% of the experts newly activated by policy $\pi_{\theta}$ for a particular rollout differ from those activated by the previous policy $\pi_{\theta_\text{old}}$. 

This effect, which is accentuated in larger MoE architectures, results in considerable oscillation of the token-level importance ratios $w_{i,t}(\theta) = \frac{ \pi_{\theta} (y_{i,t} | x, y_{i,<t}) }{ \pi_{\theta_\text{old}} (y_{i,t} | x,y_{i,<t})}$. As outlined in \S~\ref{sec:motivation} and~\ref{subsec:gradient}, this instability undermines the reliability of these weights and, consequently, interferes with the RL training's capacity to reach convergence.

\paragraph{Our Previous Approach} In our earlier work, we utilized the \textbf{Routing Replay} strategy during training to address this problem. The method involves recording the active experts determined by $\pi_{\theta_\text{old}}$ and subsequently “replaying” these same routing configurations within $\pi_{\theta}$ when calculating the importance ratios $w_{i,t}(\theta) = \frac{ \pi_{\theta} (y_{i,t} | x, y_{i,<t}) }{ \pi_{\theta_\text{old}} (y_{i,t} | x,y_{i,<t})}$. Through this approach, both $\pi_{\theta} (y_{i,t} | x, y_{i,<t})$ and $\pi_{\theta_\text{old}} (y_{i,t} | x,y_{i,<t})$ for a token $y_{i,t}$ utilize identical expert subnetworks, thereby reinstating the reliability of the token-level importance ratio and facilitating the consistent optimization of the same activated network during gradient steps. As illustrated in Figure~\ref{fig:routing_replay}, Routing Replay is vital for the stable convergence of GRPO training in MoE frameworks.

\paragraph{Benefit of GSPO} While Routing Replay permits successful convergence during the GRPO training of MoE models, it introduces extra memory usage and communication costs by repeatedly employing the same routing patterns, thereby restricting the actual expressiveness and capacity of the MoE architecture. GSPO, illustrated in Figure~\ref{fig:results}, discards reliance on Routing Replay, enabling conventional calculation of importance ratios $s_i(\theta)$ and achieving reliable convergence and optimization. The essential advantage is that GSPO considers only the likelihood of entire sequences (i.e., $\pi_{\theta} (y_i | x)$), remaining unaffected by fluctuations at the token level (i.e., $\pi_{\theta} (y_{i,t} | x, y_{i,<t})$). As the MoE model consistently upholds its language modeling abilities, sequence-level likelihood remains relatively stable. Therefore, GSPO inherently addresses volatility in expert activation for MoE models, removing the necessity for elaborate techniques such as Routing Replay. This improvement leads to a more straightforward and robust training process and allows unrestricted utilization of model capacity.

\subsection{Benefit of GSPO for RL Infrastructure}

Due to variations in numerical precision between training frameworks such as Megatron and inference frameworks like SGLang and vLLM, standard practice involves recalculating the likelihoods of generated responses using the training engine under the previous policy $\pi_{\theta_\text{old}}$. However, GSPO relies on sequence-level likelihoods for optimization rather than more granular token-level likelihoods; sequence-level likelihoods are inherently more robust against precision inconsistencies across engines. Consequently, GSPO enables optimization to proceed with the likelihood estimates provided directly by the inference engine, eliminating the necessity for recomputation within the training environment. This advantage is particularly pronounced in contexts such as partial rollouts, multi-turn reinforcement learning, and frameworks where training and inference processes are decoupled.

\section{Conclusion}

Group Sequence Policy Optimization (GSPO) is introduced as a novel reinforcement learning approach for large language model training. GSPO utilizes importance sampling by constructing importance ratios from sequence probabilities, then applies sequence-level operations such as clipping, reward assignment, and optimization. Compared to GRPO, GSPO significantly enhances stability, efficiency, and overall training outcomes, proving especially advantageous in extensive RL applications involving MoE models. These advancements underpin the remarkable progress seen in the newest Qwen3 models. GSPO thus serves as a robust and scalable foundation, paving the way for further expansion of RL and anticipatory breakthroughs in artificial intelligence.

\end{document}